We prove that invariant distributions on the rational normal spaces
classified in Section~\ref{sec:unitary-geometry} are fixed by the
corresponding corrected involutions. For regular rank-zero orbits,
we either reduce to the linear two-leg theorem or use Hermitian
Gram fibers. The scalar rank-zero case is treated by partial Fourier
transform and homogeneity. Together with generalized descent,
these arguments prove Theorem~\ref{thm:main-unitary}.

\subsection{Elementary descendants and simultaneous conjugation}
\label{m2:sec:elementary}

We first prove the elementary vanishing statements for the
representations in \cref{m2:nor:classification}. An odd
operator is always the corrected operator of the relevant
stabilizer. Multiplication of it by an even group element does not
change its action on even-invariant distributions.

We apply \cref{m2:elem:split,m2:elem:norm,m2:elem:products}
through the rational images of these stabilizers.

\subsubsection{Rational root maps on normal spaces}

For a regular toral dense orbit, let $T$ be its rational
stabilizer and $T_{\mathrm{ad}}$ its effective one-dimensional
torus. The root action computed in \cref{m2:nor:classification}
fits into
\begin{equation}\label{m2:elem:rootsequence}
 1\longrightarrow\Res_{E/F}\Gm\longrightarrow T
 \xrightarrow{\chi}T_{\mathrm{ad}}\longrightarrow1.
\end{equation}
By \cref{m2:rat:adjoint-torus}, $\chi$ is the eigenvalue ratio,
with scalar kernel, and is surjective on rational points;
$T_{\mathrm{ad}}$ is $\Gm$ in the split case and $K_T^1$ otherwise.
For the rank-one families, \cref{m2:nor:rankone-rational}
proves that the characters \eqref{m2:nor:rational-characters}
give surjections $T_3(F)\to(E^1)^2$ and $T_2(F)\to E^1$.
Thus the effective actions are the opposite-weight and
norm-one actions of \cref{m2:elem:split,m2:elem:norm,m2:elem:products}.

\begin{proposition}\label{m2:elem:vanishing}
The spaces of sign-equivariant distributions on the reduced normal
spaces vanish for the dense absolutely
irreducible, dense split toral, dense anisotropic toral,
rank-one decomposable, and rank-one stable-system families.
\end{proposition}

\begin{proof}
The calculations in \cref{m2:nor:classification} identify the
following extended representations on the reduced normal spaces.
In the absolutely irreducible case, the normal space is $F^5$ and
the odd operator acts trivially, so sign equivariance forces zero.
The split toral case is $F^4\oplus F^2$ with the last two
coordinates of opposite weights and swapped by the odd map.
The anisotropic toral case is $F^4\oplus K_T$, with
$K_T^1$ acting on the last coordinate and conjugation odd.
The rank-one cases are $F^3\oplus E^2$ and $F^4\oplus E$,
with the coordinate norm actions and conjugations just described.
The indicated copies of $F$ are fixed by the entire extended
stabilizer. The fixed coordinates are passive by \cref{m2:tools:passive}.
The rational root surjections then permit the application of
\cref{m2:elem:split,m2:elem:norm,m2:elem:products} to the remaining
coordinates, proving all five assertions.
\end{proof}

\subsubsection{Simultaneous-conjugation vanishing}\label{m2:sec:simconj}

We use the rational factorization, the coordinates $\Psi$, and the
matrices $J_0,R$ of \cref{m2:red:unitary-pair}, together with
\eqref{m2:simconj:transpose}.

\begin{theorem}\label{m2:simconj:vanishing}
Every $G_2(F)$-conjugation-invariant distribution on $M_2(E)$
is invariant under ordinary transpose.
\end{theorem}

\begin{proof}
Let $T'=\Psi_*T$. By \cref{m2:simconj:factorization}, it is invariant
under simultaneous $\GL_2(F)$-conjugation on $M_2(F)^2$.
The common matrix-pair result \cref{m1:lem:matrix-pair-transpose}
therefore fixes $T'$ under componentwise transpose. The odd
coordinate formula \eqref{m2:simconj:transpose} adds conjugation by
$R\in\GL_2(F)$, which also fixes $T'$. Pulling back through $\Psi$
proves the assertion.
\end{proof}

\subsection{Regular unitary two-leg representations}\label{m2:leg:section}\label{m2:sec:twoleg}

Let $D\in G_2(F)$ be regular semisimple, and put
\[
 T=Z_{G_2}(D),\qquad S=G_2\times_{\lambda}T,
 \qquad \mathcal L=\Mat_2(E).
\]
The even action is
\begin{equation}\label{m2:leg:action}
 (v,h)\cdot Q=vQh^{-1}.
\end{equation}
Throughout the regular-descendant argument the odd operator is specified by matrices
$c_U,c_W\in G_2(F)$ satisfying
\begin{equation}\label{m2:leg:transport-condition}
 \lambda(c_U)=\lambda(c_W),\qquad c_WD^tc_W^{-1}=D,
\end{equation}
and is
\begin{equation}\label{m2:leg:odd}
 \Theta(Q)=c_U\kappa(Q)c_W^{-1},
 \qquad \kappa(Q)=-J\bar QJ.
\end{equation}
These are precisely the conditions furnished by the regular rank-zero
normal representation in \cref{m2:nor:zero-normal,m2:nor:zero-odd}.
We first check the extended action, then prove its distribution theorem.

\begin{lemma}\label{m2:leg:extended-action}
The operator $\Theta$ normalizes the action of $S(F)$, and $\Theta^2$
belongs to that even action. In particular it makes sense to require a
distribution to be $S(F)$-invariant and $\Theta$-anti-invariant.
Moreover,
\begin{equation}\label{m2:leg:gram-odd}
 \Theta(Q)^*\Theta(Q)=c_W(Q^*Q)^tc_W^{-1}.
\end{equation}
\end{lemma}

\begin{proof}
For $g\in G_2(F)$, the unitary equation gives
$J\bar gJ=\lambda(g)g^{-t}$. Hence on an equal-multiplier pair,
\[
 \kappa(vQh^{-1})=v^{-t}\kappa(Q)h^t.
\]
Every $h\in T(F)$ belongs to $E[D]$, because the centralizer of a
regular two-by-two matrix is $E[D]$.
Condition \eqref{m2:leg:transport-condition} consequently gives
\begin{equation}\label{m2:leg:transpose-centralizer}
 c_Wh^tc_W^{-1}=h\quad(h\in T(F)),\qquad c_WT^tc_W^{-1}=T.
\end{equation}
Thus $\Theta$ conjugates $(v,h)$ to
$(c_Uv^{-t}c_U^{-1},c_Wh^{-t}c_W^{-1})\in S(F)$.
Their multipliers are both $\lambda(v)^{-1}$.
Since $\kappa^2=1$, a second application yields
\begin{equation}\label{m2:leg:odd-square}
 \Theta^2(Q)=(c_Uc_U^{-t})Q(c_Wc_W^{-t})^{-1}.
\end{equation}
Transposing the second condition in \eqref{m2:leg:transport-condition}
shows that $c_Wc_W^{-t}$ commutes with $D$. Both factors in
\eqref{m2:leg:odd-square} have multiplier one, proving the assertion.

For the Gram formula, $\kappa(Q)^*=-Q^t$, so
\[
 \kappa(Q)^*\kappa(Q)=Q^tJ\bar QJ=(Q^*Q)^t.
\]
Using $c_U^*c_U=\lambda(c_U)I$ and
$(c_W^{-1})^*=\lambda(c_W)^{-1}c_W$ now proves
\eqref{m2:leg:gram-odd}. The equality of multipliers is essential in
this cancellation.
\end{proof}

We write $\Ddist(\mathcal L)^{S,\Theta=-1}$ for the distributions
that are $S(F)$-invariant and $\Theta$-anti-invariant.

\begin{theorem}\label{m2:leg:regular}
For every regular semisimple $D\in G_2(F)$ and every pair
$c_U,c_W$ satisfying \eqref{m2:leg:transport-condition},
\[
 \Ddist(\Mat_2(E))^{S,\Theta=-1}=0.
\]
Equivalently, every distribution invariant under \eqref{m2:leg:action}
is fixed by \eqref{m2:leg:odd}.
\end{theorem}

The proof divides according to whether the adjoint torus $T/Z(G_2)$
is split or anisotropic.

\subsubsection{Rational factorization and the effective action}

As in \cref{m2:red:unitary-pair}, we have $\bar\delta=-\delta$,
$d=\delta^2$, $P=\diag(1,\delta)$ and $R=P^{-2}$.
By \cref{m2:simconj:factorization}, write
\[
 D=z_DPg_DP^{-1},\qquad z_D\in E^\times,
 \quad g_D\in\GL_2(F).
\]
The matrix $g_D$ is regular semisimple. Set
\[
 T_0(F)=Z_{\GL_2(F)}(g_D)=F[g_D]^\times.
\]
If $h=zPgP^{-1}$, then $hD=Dh$ is equivalent to $gg_D=g_Dg$.
It follows directly that
\begin{equation}\label{m2:leg:rational-torus}
 T(F)=E^\times PT_0(F)P^{-1}.
\end{equation}
The common scalar overlap is exactly $F^\times$:
$E^\times I\cap P\GL_2(F)P^{-1}=F^\times I$.
In particular the quadratic \emph{étale} $F$-algebra $F[g_D]$
determines the split or anisotropic type of the adjoint torus.

The kernel of \eqref{m2:leg:action} is the diagonal scalar group
\begin{equation}\label{m2:leg:kernel}
 \{(zI,zI):z\in E^\times\}.
\end{equation}
Indeed an element in the kernel has $v=h$ by evaluation at $Q=I$;
that matrix then commutes with every $Q\in\Mat_2(E)$ and is scalar.
Conversely, every pair $(zI,zI)$ in \eqref{m2:leg:kernel} has equal multipliers and
acts trivially.

Write $v=zPgP^{-1}$ and $h=wPtP^{-1}$, with
$g\in\GL_2(F)$ and $t\in T_0(F)$. Put $r=z/w$.
Since $\lambda(PgP^{-1})=\det g$, the multiplier condition becomes
\begin{equation}\label{m2:leg:effective-triples}
 N_{E/F}(r)\det g=\det t.
\end{equation}
Changing the two scalar lifts produces exactly
\[
 (r,g,t)\sim((b/a)r,ag,bt),\qquad a,b\in F^\times.
\]
If $P^{-1}QP=A+\delta B$ and $r=a_0+\delta b_0$, the action is
\begin{equation}\label{m2:leg:mixing}
 (A,B)\longmapsto
 \bigl(g(a_0A+db_0B)t^{-1},\ g(b_0A+a_0B)t^{-1}\bigr).
\end{equation}
This follows by multiplying $r g(A+\delta B)t^{-1}$.
The two rational matrix components are therefore mixed by a general
even element. We use a different linear map in the split case.

\subsubsection{Reduction for split adjoint tori}

Suppose $F[g_D]\simeq F\times F$. A change of basis by
$b=PfP^{-1}\in G_2(F)$ diagonalizes $D$. Such a change is compatible
with all the specified data: replace $D,Q$ by $bDb^{-1},bQb^{-1}$,
replace the even factors by their conjugates, and replace
\[
 c_U,c_W\quad\hbox{by}\quad bc_Ub^t,bc_Wb^t.
\]
Substitution in \eqref{m2:leg:odd} gives exactly the new odd operator.
Both new multipliers are multiplied by $\lambda(b)^2$, and the
transpose-transporter condition is preserved. We may thus assume
that $D$ and $T$ are diagonal in the hyperbolic basis with Gram matrix
$J$. An element $h=\diag(h_1,h_2)$ of $T(F)$ satisfies
$\bar h_1h_2\in F^\times$.

\begin{lemma}\label{m2:leg:split-effective}
The map
\begin{equation}\label{m2:leg:split-group-map}
 (g,s)\longmapsto
 \left(PgP^{-1},\ h_{g,s}\right),\qquad
 h_{g,s}=\diag\left(s,\frac{\det g}{\bar s}\right),
\end{equation}
identifies the effective action of $S(F)$ with that of
$(\GL_2(F)\times E^\times)/\diag F^\times$.
\end{lemma}

\begin{proof}
Both matrices in \eqref{m2:leg:split-group-map} have multiplier $\det g$,
and the formula is multiplicative. For surjectivity on effective
actions, factor an arbitrary $v=zPgP^{-1}$ and divide the pair
$(v,h)$ by $(zI,zI)$. The remaining right matrix is diagonal, of
multiplier $\det g$, so its second diagonal entry is
$\det g/\bar s$ if its first entry is $s$.
Finally the pair acts trivially precisely when
$PgP^{-1}=h_{g,s}$ is scalar. This means $g=aI$, $s=a$,
$a\in F^\times$. This proves the asserted kernel as well.
\end{proof}

Regard $E$ as an $F$-vector space with ordered basis $(1,\delta)$, and
let $m_s$ be multiplication by $s$ on this space. Define
\[
 \mathcal T(x)(u)=
 \begin{pmatrix}\Tr_{E/F}(x_1u)\\\Tr_{E/F}(x_2u)\end{pmatrix},
 \qquad x=(x_1,x_2)^t\in E^2.
\]
The trace pairing is nondegenerate. More explicitly, if
$x_i=x_{i0}+\delta x_{i1}$, then
\begin{equation}\label{m2:leg:trace-matrix}
 \mathcal T(x)=2
 \begin{pmatrix}x_{10}&d x_{11}\\x_{20}&d x_{21}\end{pmatrix}.
\end{equation}
For a two-by-two matrix over $F$ write
$Y^\#=\tr(Y)I-Y$; this is the adjugate, and is an $F$-linear
involution. Write $P^{-1}QP=(u,v)$ by columns and define
\begin{equation}\label{m2:leg:linear-map}
 \mathcal I(Q)=(X,Y),\qquad X=\mathcal T(u),
 \quad Y^\#=\mathcal T(v).
\end{equation}
It is an $F$-linear homeomorphism
\[
 \Mat_2(E)\longrightarrow
 \Hom_F(E,F^2)\oplus\Hom_F(F^2,E).
\]

\begin{lemma}\label{m2:leg:split-intertwining}
Under \eqref{m2:leg:split-group-map} and \eqref{m2:leg:linear-map}, the
even action becomes
\begin{equation}\label{m2:leg:model-action}
 (g,s)(X,Y)=(gXm_s^{-1},m_sYg^{-1}).
\end{equation}
Let $S_E=\diag(2,2d)$, the Gram matrix of
$\Tr_{E/F}(uu')$, and put
\begin{equation}\label{m2:leg:model-odd}
 \tau_E(X,Y)=(Y^tS_E,S_E^{-1}X^t).
\end{equation}
Then the odd intertwining formula is
\begin{equation}\label{m2:leg:odd-intertwining}
 \mathcal I\kappa\mathcal I^{-1}
   =(R,m_{2\delta})\cdot\tau_E.
\end{equation}
Here the pair on the right is an element of the even group in
\eqref{m2:leg:model-action}.
\end{lemma}

\begin{proof}
The two $E$-columns transform by
\[
 u\longmapsto gu s^{-1},\qquad
 v\longmapsto(\det g)^{-1}gv\bar s.
\]
We have $\mathcal T(xs)=\mathcal T(x)m_s$. Also adjugation reverses
products and
\[
 (g^{-1})^\#=(\det g)^{-1}g,
 \qquad m_s^\#=N_{E/F}(s)m_s^{-1}=m_{\bar s}.
\]
Consequently
\[
 (m_sYg^{-1})^\#=(\det g)^{-1}gY^\#m_{\bar s},
\]
which proves \eqref{m2:leg:model-action} for both legs.

For the odd formula set
$L=\left(\begin{smallmatrix}0&-\delta\\\delta^{-1}&0\end{smallmatrix}\right)$.
Direct substitution of $Q=P(u,v)P^{-1}$ gives
\[
 P^{-1}\kappa(Q)P=L\overline{(u,v)}L.
\]
Its columns are
\[
 \begin{pmatrix}-\bar v_2\\d^{-1}\bar v_1\end{pmatrix},
 \qquad
 \begin{pmatrix}d\bar u_2\\-\bar u_1\end{pmatrix}.
\]
Using \eqref{m2:leg:trace-matrix} and
$Y^\#=J_0Y^tJ_0^{-1}$ yields
\begin{equation}\label{m2:leg:odd-coordinates}
 \mathcal I\kappa\mathcal I^{-1}(X,Y)
       =(RY^tJ,JX^tR^{-1}).
\end{equation}
Since
\[
 m_{2\delta}=\begin{pmatrix}0&2d\\2&0\end{pmatrix}=JS_E,
 \qquad S_Em_{2\delta}^{-1}=J,
\]
applying $(R,m_{2\delta})$ to \eqref{m2:leg:model-odd} gives exactly
\eqref{m2:leg:odd-coordinates}.
\end{proof}

We apply the linear two-leg theorem \cref{m1:thm:two-leg}, using
its trace-form adjoint \eqref{m1:eq:two-leg-adjoint}.

In our application
$V=F^2$, $W=E$, $T'=E^\times$ acts by multiplication, and
$S_W=S_E$. The trace identity
$\Tr((su)u')=\Tr(u(su'))$ gives
$S_Em_s=m_s^tS_E$. Thus this is the nonsplit case of that theorem,
even though the unitary adjoint torus considered here is split.
The scalar subgroup $(aI,aI)$ acts trivially; passage to its
effective quotient changes neither the distribution space nor the
sign character. The identification is on $F$-points and on the entire
linear representation.

For singular legs, the nonsplit proof is given in \cref{m1:subsec:two-leg-nonsplit}, using
\cref{m1:lem:matrix-boundary,m1:lem:rank-one-factorization,m1:lem:common-U11}. It localizes by the complete composition invariants
of $YX$, then treats ranks two, one, and zero of that composition.
For rank one, its factorization lemma includes singular individual
legs. For zero composition, it first treats the open locus
$\rk X=\rk Y=1$, and only then the closed union of matrix axes.
On a rank-one matrix axis, descent over $\mathbf P(V)$ leaves the
upper triangular group
\[
 P_V=\left\{\begin{pmatrix}\alpha&b\\0&\beta\end{pmatrix}\right\}
\]
with required Frobenius character $|\beta/\alpha|$; the element
$\diag(1,\beta)$ acts trivially on the resulting geometric fiber.
This is the modular obstruction on that stratum. On the two rank-two
axes the finite-part distributions have opposite logarithmic defects
and the same nonzero residue at the origin. 

We can finish the split case of \cref{m2:leg:regular}. For diagonal
regular $D$, \eqref{m2:leg:transport-condition} says that $c_W$ commutes
with $D$, so $c_W\in T(F)$ and $(c_U,c_W)\in S(F)$.
Thus $\Theta$ differs from $\kappa$ by an even action. By
\cref{m2:leg:split-intertwining}, $\kappa$ in turn differs from
$\tau_E$ by the explicit even pair $(R,m_{2\delta})$.
Thus \cref{m1:thm:two-leg} makes every even-invariant distribution
$\Theta$-invariant.

\subsubsection{Anisotropic tori and the self-adjoint Gram space}

Assume now that $T/Z(G_2)$ is anisotropic. Let $W=E^2$ be the right
Hermitian space, let $M=E[D]\subset\End_E(W)$, and denote the
Hermitian adjoint on $M$ by $*$. Since
$D^*=\lambda(D)D^{-1}$, this algebra is stable under $*$.
Put $L=M^{*=1}$. There are two forms to consider:
\begin{align}
 &M=E\times E,\qquad
   (r_1,r_2)^*=(\bar r_1,\bar r_2);\label{m2:leg:an-product}\\
 &M=EL\text{ a biquadratic field},\qquad
   \langle *\rangle=\Gal(M/L).\label{m2:leg:an-field}
\end{align}
These are the anisotropic cases of \cref{m2:rat:torus-cohomology}.
Let $\iota$ denote the nontrivial $E$-automorphism of $M$.
If $K_T$ is the splitting field of the adjoint torus, then
\begin{equation}\label{m2:leg:splitting-field}
 K_T=E\quad\hbox{in \eqref{m2:leg:an-product}},\qquad
 K_T=M^{*\iota}\quad\hbox{in \eqref{m2:leg:an-field}}.
\end{equation}
These identifications are those of \cref{m2:rat:adjoint-torus}.
In both cases
\[
 T(F)=\{r\in M^\times:r^*r\in F^\times\}.
\]

Since $D$ is cyclic, a cyclic vector identifies the rank-one
$M$-module $W$ with $M$. The trace-form calculation in
\cref{m2:rat:torus-cohomology} writes its Hermitian form as
\begin{equation}\label{m2:leg:trace-hermitian}
 \langle u,v\rangle=\Tr_{M/E}(c uv^*),
 \qquad c\in(M^{*=1})^\times.
\end{equation}
These coordinates carry the given hyperbolic form on $W$.

Put
\[
 \mathcal H=\{A\in\End_E(W):A^*=A\},\qquad
 \Gamma(Q)=Q^*Q.
\]
The action of $S$ on Gram matrices is
\begin{equation}\label{m2:leg:gram-even}
 \Gamma(vQh^{-1})=h\Gamma(Q)h^{-1}.
\end{equation}
For example this follows from
$(h^{-1})^*=\lambda(h)^{-1}h$ and the equality
$\lambda(v)=\lambda(h)$. We compute the rational $T$-orbits of
$\mathcal H$ before analyzing the Gram factorization fibers.

\begin{lemma}\label{m2:leg:root-image}
There is a $T$-stable decomposition over $F$
\[
 \mathcal H=C\oplus W_{\mathrm{rt}},\qquad
 \dim_F C=2,\qquad W_{\mathrm{rt}}\simeq K_T,
\]
where $C$ is pointwise fixed and $K_T^1$ acts on the
root plane by multiplication. The rational root maps have kernel
$E^\times$ and are surjective.
\end{lemma}

\begin{proof}
In \eqref{m2:leg:an-product}, write $c=(c_1,c_2)$ with $c_i\in F^\times$.
Relative to the two idempotents, the self-adjoint matrices are
\begin{equation}\label{m2:leg:product-selfadjoint}
 A=\begin{pmatrix}a&b\\(c_1/c_2)\bar b&e\end{pmatrix},
 \qquad a,e\in F,\quad b\in E.
\end{equation}
The fixed coordinates are $a$ and $e$; the root coordinate is $b$.
Conjugation by $(r_1,r_2)\in T$ multiplies $b$ by $r_1/r_2$.
This is the root ratio in \cref{m2:rat:adjoint-torus}.

In \eqref{m2:leg:an-field}, $c\in L^\times$ and every $E$-endomorphism
of $M$ has a unique expression $m_x+m_y\iota$, $x,y\in M$.
The two summands are linearly independent: otherwise $\iota$ would
be multiplication by an element of $M$, contrary to its action on $1$
and on an element outside $E$. Dimensions then give the assertion.
Taking the adjoint in \eqref{m2:leg:trace-hermitian} shows
\[
 (m_y\iota)^*=m_{c^{-1}\iota(c)\iota(y^*)}\iota.
\]
Consequently the self-adjoint matrices are precisely
\begin{equation}\label{m2:leg:field-selfadjoint}
 A=m_x+m_{c^{-1}z}\iota,\qquad x\in L,\quad z\in K_T.
\end{equation}
Here $y=c^{-1}z$ satisfies the adjoint identity above precisely
when $z=\iota(z^*)$. Multiplication operators $m_x$ give $C$, and
the $z$-plane is the root plane.

Conjugation by $r\in T$ multiplies $z$ by $r/\iota(r)$.
This is again the root ratio of \cref{m2:rat:adjoint-torus},
which gives its rational surjectivity and scalar kernel in both cases.
\end{proof}

Choose any $F$-basis of $C$ and write a self-adjoint matrix as
$(a_1,a_2,w)\in F^2\oplus K_T$. Define
\begin{equation}\label{m2:leg:an-invariants}
 \pi(Q)=\bigl(a_1(\Gamma(Q)),a_2(\Gamma(Q)),
                     N_{K_T/F}(w(\Gamma(Q)))\bigr)\in F^3.
\end{equation}

\begin{lemma}\label{m2:leg:gram-orbits}
Every nonempty fiber of $(a_1,a_2,N(w))$ on $\mathcal H$ is one
rational $T(F)$-orbit. Its points all have the same rank.
The map \eqref{m2:leg:an-invariants} is invariant under both $S$ and
$\Theta$.
\end{lemma}

\begin{proof}
A nonzero norm fiber in the field $K_T$ is one $K_T^1$-orbit:
if $N(w')=N(w)\ne0$, then $w'/w\in K_T^1$.
The zero norm fiber consists only of zero. Rational surjectivity in
\cref{m2:leg:root-image} therefore gives the orbit assertion.
In product coordinates \eqref{m2:leg:product-selfadjoint},
\[
 \det_E A=ae-(c_1/c_2)N_{E/F}(b).
\]
In field coordinates \eqref{m2:leg:field-selfadjoint},
\begin{equation}\label{m2:leg:field-determinant}
 \det_E A=N_{L/F}(x)-N_{L/F}(c)^{-1}N_{K_T/F}(z).
\end{equation}
For the latter identity, after extending scalars to split $M/E$,
$m_x+m_y\iota$ has matrix
$\left(\begin{smallmatrix}x&y\\\iota(y)&\iota(x)\end{smallmatrix}\right)$;
its determinant is $x\iota(x)-y\iota(y)$, giving the formula over $F$.
These formulas determine the determinant on every invariant fiber.
They also distinguish the zero matrix: in the product case all its
coordinates vanish; in the field case $z=0,x\ne0$ gives an invertible
multiplication operator. Rank is consequently constant.

Even invariance follows from \eqref{m2:leg:gram-even}.
For odd invariance put $\vartheta(A)=c_WA^tc_W^{-1}$.
Transpose preserves $\mathcal H$: if $A^*=A$, then
$A^t=J\bar AJ$ and $(A^t)^*=J\bar AJ=A^t$.
Conjugation by a unitary similitude also preserves self-adjointness,
so $\vartheta$ preserves $\mathcal H$.
It fixes $C$ pointwise: every element of $C$ is in $E[D]$, and
$c_Wf(D)^tc_W^{-1}=f(D)$.
Furthermore \eqref{m2:leg:transpose-centralizer} gives
\[
 \vartheta(hAh^{-1})=h^{-1}\vartheta(A)h.
\]
Thus on the nontrivial root summand it intertwines multiplication
by $q$ with multiplication by $q^{-1}$, $q\in K_T^1$.
It follows that it has the form $w\mapsto\eta\bar w$:
composing with field conjugation gives an $F$-linear map commuting
with $K_T^1$, hence with the algebra $K_T$ which that infinite group
spans. Such a map is multiplication by an element $\eta\in K_T^\times$.
Since $\vartheta$ preserves determinants and fixes the $C$-coordinates,
the determinant formulas force $N(\eta)=1$.
It therefore fixes $N(w)$, proving odd invariance of \eqref{m2:leg:an-invariants}.
\end{proof}

\subsubsection{Nonzero Gram fibers}

For a represented $A\in\mathcal H$ let
\[
 \mathcal F_A=\{Q:Q^*Q=A\},\qquad T_A=Z_T(A),
 \qquad S_A=G_2\times_\lambda T_A.
\]
The base $T$-orbit of $A$ is homogeneous for $S$, with stabilizer $S_A$,
and its Gram fiber is $\mathcal F_A$.
By \cref{m2:leg:root-image}, $T_A$ is either $T$ or the scalar torus
$E^\times$; in particular $S_A$ is reductive and unimodular.
Thus Frobenius descent over this base has relative modulus one.

\begin{lemma}\label{m2:leg:gram-factorization}
If $A$ has rank two, $U_{1,1}(F)$ acts simply transitively on
$\mathcal F_A$, and the $S$-stabilizer of $Q\in\mathcal F_A$ is
\begin{equation}\label{m2:leg:rank-two-stabilizer}
 S_Q=\{(QhQ^{-1},h):h\in T_A\}.
\end{equation}
If $A$ has rank one and is nonzero, $\mathcal F_A$ is one rational
$U_{1,1}(F)$-orbit. Its $S$-stabilizer is the torus
\begin{equation}\label{m2:leg:rank-one-stabilizer}
 S_Q\simeq\{(h,s)\in T_A\times E^\times:
                                      N_{E/F}(s)=\lambda(h)\}.
\end{equation}
In both cases $\Delta_S|_{S_Q}/\Delta_{S_Q}=1$.
\end{lemma}

\begin{proof}
For rank two, $Q$ is invertible. If $Q_1^*Q_1=Q_2^*Q_2$, then
$Q_2Q_1^{-1}$ has adjoint times itself equal to $I$. Its action is
the unique unitary element taking $Q_1$ to $Q_2$.
The stabilizer equation $vQ=Qh$ forces $v=QhQ^{-1}$, and the
similitude condition for $v$ is equivalent to $hA=Ah$.
This proves \eqref{m2:leg:rank-two-stabilizer}.

For rank-one nonzero $A$, $Q$ has rank one: it cannot be invertible,
and it cannot be zero. Its image cannot be an isotropic line,
since that would give $Q^*Q=0$. Thus its image is a nondegenerate
Hermitian line. Since $\ker Q\subset\ker A$ and both are lines,
$\ker Q=\ker A$. For two factorizations of $A$ the rule
$Q_1w\mapsto Q_2w$ is therefore well-defined. It is an isometry,
because
\[
 \langle Q_iw,Q_iw'\rangle=\langle w,Aw'\rangle.
\]
To construct the unitary extension, choose $w_0$ whose image under either factorization
is nonzero and put $e_i=Q_iw_0$. Their common norm
$a=\langle w_0,Aw_0\rangle$ is nonzero. Choose a nonzero
$f_i\in(Ee_i)^\perp$ and write $b_i=\langle f_i,f_i\rangle$.
The complementary lines are nondegenerate, so $b_i\ne0$.
Let $B_i$ be the $E$-matrix with columns $e_i,f_i$ in the original
hyperbolic basis. Then
\[
 B_i^\dagger JB_i=\diag(a,b_i),\qquad
 ab_i=-N_{E/F}(\det B_i).
\]
Thus $b_1/b_2=N_{E/F}(\det B_1/\det B_2)$.
The map sending $e_1$ to $e_2$ and $f_1$ to
$(\det B_1/\det B_2)f_2$ is an isometry of both orthogonal lines.
It is therefore the required rational unitary extension on the
whole plane. This proves transitivity and supplies explicitly the
two-dimensional Hermitian Witt extension used here.

If $h$ commutes with $A$, it preserves the Hermitian form
$\langle w,Aw'\rangle$ with multiplier $\lambda(h)$.
On its nondegenerate one-dimensional quotient it therefore has an
eigenvalue of norm $\lambda(h)$. The equation $vQ=Qh$ uniquely
determines $v$ on $\im Q$ with this eigenvalue. On its orthogonal
complement $v$ is an arbitrary scalar $s$ of norm $\lambda(h)$.
This is \eqref{m2:leg:rank-one-stabilizer}; conversely each such choice
is a stabilizing similitude. The stabilizers in \eqref{m2:leg:rank-two-stabilizer} and
\eqref{m2:leg:rank-one-stabilizer} are tori, so they and $S$ are unimodular.
\end{proof}

Consider a nonempty fiber of \eqref{m2:leg:an-invariants} whose Gram
orbit is nonzero. It is one $S(F)$-orbit. Indeed, for two points first
choose $h\in T(F)$ conjugating their Gram matrices. Choose
$v_0\in G_2(F)$ of multiplier $\lambda(h)$, for example
$\diag(1,\lambda(h))$. Then $v_0Q_1h^{-1}$ and $Q_2$ have equal
Gram matrices, and \cref{m2:leg:gram-factorization} connects them by
a unitary left factor.
The preceding lemma supplies its unimodular stabilizer, so its
invariant distribution space is generated by a positive orbital
measure. The complete fiber is $\Theta$-stable by
\cref{m2:leg:gram-orbits}. The pushforward of this positive measure is
a positive scalar multiple of itself, and \eqref{m2:leg:odd-square}
forces that scalar to square to one. It is therefore one.
Consequently every nonzero complete Gram-invariant fiber has zero
sign-distribution space. 

\subsubsection{The zero Gram fiber and its nonreductive stabilizers}

It remains to prove vanishing on the entire closed fiber
\[
 \mathcal N=\Gamma^{-1}(0),\qquad
 \mathcal N^\times=\mathcal N\setminus\{0\}.
\]
For $Q\in\mathcal N$, its image is totally isotropic. The target is
a hyperbolic Hermitian plane, so every nonzero such $Q$ has rank one
and can be written
\begin{equation}\label{m2:leg:null-factorization}
 Q=e\ell,\qquad e\ne0\text{ isotropic},\quad
 \ell\in W^\vee\setminus\{0\}.
\end{equation}
Its kernel is an $E$-line. Two such maps lie in the same $S$-orbit
exactly when their kernel lines lie in the same $T$-orbit.
Necessity follows from \eqref{m2:leg:action}. For sufficiency, use $h$
to align the kernel lines. The remaining maps have proportional
nonzero functionals, and their isotropic image vectors can be
matched by a similitude of any prescribed multiplier.
Explicitly, in hyperbolic bases beginning with the respective
image vectors, an element with first eigenvalue $a\in E^\times$
and prescribed multiplier $\lambda$ has second eigenvalue
$\lambda/\bar a$. Taking $\lambda=\lambda(h)$ gives the required
left factor, with no norm restriction on $a$.

Under the identification $W\simeq M$ of
\eqref{m2:leg:trace-hermitian}, the open locus of projective lines
with an invertible representative is
$M^\times/E^\times$. Its $T$-orbit space is
\begin{equation}\label{m2:leg:null-quotient}
 \Xi=M^\times/(E^\times T(F))=M^\times/T(F),
\end{equation}
since $E$-scalars already lie in $T$.
For a field $M$, every nonzero vector is invertible, so this is the
whole projective-line quotient. For the product algebra the missing
lines are its two idempotent eigenlines, which we treat below.

The group $T(F)$ is closed in $M^\times$, and its projective action
factors through the compact group $T(F)/E^\times\simeq K_T^1$.
The quotient of the compact projective $E$-line by this compact action
is Hausdorff and totally disconnected. For the latter assertion,
given an orbit and an invariant open neighborhood of it, compactness
of the orbit gives a clopen set $C$ containing the entire orbit and
contained in that neighborhood. Compactness and continuity show
that $C$ has open stabilizer in the compact acting group, so its
translates form a finite family. The intersection of this family
still contains the entire orbit and is invariant and clopen.
The images of these neighborhoods provide
a clopen basis in the quotient. The image of the invariant
open invertible locus is open, and identifies with \eqref{m2:leg:null-quotient}.
In particular $\Xi$ is a separated locally compact totally
disconnected space. The kernel-line map on the rank-one locus is
continuous. Its composition with \eqref{m2:leg:null-quotient} is
$S$-invariant, with each fiber a whole rational $S$-orbit, by the
preceding orbit calculation.

\begin{lemma}\label{m2:leg:null-stabilizer}
Let $Q=e\ell\in\mathcal N^\times$, let $T_{\ker\ell}$ be the
stabilizer of its kernel line, and define $a(h)$ by
$\ell h=a(h)\ell$. In a hyperbolic basis beginning with $e$ its
stabilizer consists of
\begin{equation}\label{m2:leg:null-matrices}
 h\in T_{\ker\ell},\qquad
 v=\begin{pmatrix}
 a(h)&a(h)u\\0&\lambda(h)/\overline{a(h)}
 \end{pmatrix},\quad u\in E^-.
\end{equation}
It has the form $E^-\rtimes T_{\ker\ell}$, and
\begin{equation}\label{m2:leg:null-modulus}
 \Delta_{S_Q}(h)=
       \left|N_{E/F}(a(h))/\lambda(h)\right|_F.
\end{equation}
For every kernel line in the present anisotropic cases this modulus
equals one.
\end{lemma}

\begin{proof}
The stabilizer equation is $v e\ell=e\ell h$, so it requires $h$
to fix the kernel and $ve=a(h)e$. The unitary-similitude equation
with multiplier $\lambda(h)$ then forces exactly
\eqref{m2:leg:null-matrices}; its upper-right condition is
$u+\bar u=0$. Conjugation by the diagonal part of $v$ acts on
the $F$-line $E^-$ by
\[
 u\longmapsto\frac{N(a(h))}{\lambda(h)}u.
\]
The torus has modulus one, and this determinant on the unipotent
radical gives \eqref{m2:leg:null-modulus}, with the convention of
\cref{m2:tools:frobenius}.

For an invertible line $Ew$, the relation $hw=aw$ forces $h=a\in E^\times$,
so $T_{Ew}=E^\times$. In this case $N(a(h))=\lambda(h)$.
For either exceptional product eigenline, its stabilizer is all $T$;
$a(h)$ is one of the two components of $h$, and
\eqref{m2:leg:an-product} gives $N(h_1)=N(h_2)=\lambda(h)$.
Thus the quotient is one in every case under consideration.
\end{proof}

We next compute the odd action on the rational orbits, using
\eqref{m2:leg:transport-condition}.

\begin{lemma}\label{m2:leg:null-odd}
The operator $\Theta$ preserves each $S(F)$-orbit in $\mathcal N^\times$.
\end{lemma}

\begin{proof}
On kernel lines the induced anti-linear map is
\[
 j(w)=c_WJ\bar w.
\]
Indeed $\ker\kappa(Q)=J\overline{\ker Q}$ and the right factor
$c_W^{-1}$ in \eqref{m2:leg:odd} then multiplies this kernel by $c_W$.
Put $\mu=\lambda(D)$. The unitary equation and the transporter give
\[
 jDj^{-1}=c_WJ\bar DJc_W^{-1}
          =\mu c_WD^{-t}c_W^{-1}
          =\mu D^{-1}=D^*.
\]
Since $j$ conjugates $E$-scalars by the nontrivial field automorphism,
it follows that $j(rw)=r^*j(w)$ for every $r\in M$.
Writing $W$ as $M$ and putting $a=j(1)$ therefore gives
\begin{equation}\label{m2:leg:null-antilinear}
 j(w)=aw^*,\qquad a\in M^\times.
\end{equation}
Invertibility of $a$ follows from invertibility of $j$.

The map $J\bar{\phantom w}$ is anti-unitary for the original form
with matrix $J$. Thus $j$ is an anti-similitude of multiplier
$\lambda(c_W)$: $\langle j(u),j(v)\rangle
=\lambda(c_W)\overline{\langle u,v\rangle}$.
In the trace coordinates \eqref{m2:leg:trace-hermitian}, this says
\[
 \Tr_{M/E}(c a^*a u^*v)
       =\lambda(c_W)\Tr_{M/E}(c u^*v)\qquad(u,v\in M).
\]
Nondegeneracy of the trace pairing and $c\in M^\times$ imply
$a^*a=\lambda(c_W)\in F^\times$, so $a\in T(F)$.
For invertible $w$,
\[
 (w^*/w)^*(w^*/w)=1,
\]
and hence $w^*/w\in T(F)$. Formula
\eqref{m2:leg:null-antilinear} therefore fixes the class of $w$ in
$\Xi$. In the componentwise product case $*$ fixes each idempotent;
multiplication by $a$ also preserves each eigenline. Those two
exceptional projective lines are consequently fixed individually.
The kernel-orbit classification proves the lemma.
\end{proof}

Every orbit just considered has the unimodular stabilizer
of \cref{m2:leg:null-stabilizer}. By Frobenius descent its invariant
distribution space is one-dimensional, generated by a positive
orbital measure. The odd operator preserves its orbit by
\cref{m2:leg:null-odd}; its positive scaling constant squares to one
by \eqref{m2:leg:odd-square}. Thus that measure is fixed.

By \cref{m2:tools:localization}, vanishing on the open
invertible-kernel locus in $\mathcal N^\times$ reduces to the complete
fibers of the map into $\Xi$. Every such fiber is an orbit already
shown to have zero sign space, so vanishing holds on the whole open locus. In the product case its
complement in $\mathcal N^\times$ is the union of the two closed
exceptional orbits. Each is odd-stable and has the same vanishing;
\cref{m2:tools:support} completes the passage to
$\mathcal N^\times$. There is no exceptional locus in the field case.
It remains to pass from $\mathcal N^\times$ to $\mathcal N$.
Any remaining sign distribution is supported at zero. Such a
distribution on this non-Archimedean space is a scalar multiple of
$\delta_0$, which is fixed by $\Theta$. Therefore
\begin{equation}\label{m2:leg:null-vanishing}
 \Ddist(\mathcal N)^{S,\Theta=-1}=0.
\end{equation}

\subsubsection{Conclusion of the regular theorem}

For the anisotropic case, every complete fiber of
\eqref{m2:leg:an-invariants} is empty, a single orbit with nonzero Gram
matrix, or the zero Gram fiber. The first case contributes nothing,
the second was treated by \cref{m2:leg:gram-factorization} and the
positive-measure argument, and the third is
\eqref{m2:leg:null-vanishing}. The map is continuous and invariant
under the whole extended action by \cref{m2:leg:gram-orbits}.
Bernstein localization therefore proves \cref{m2:leg:regular} in both
anisotropic forms. The split case was proved using
\cref{m2:leg:split-intertwining,m1:thm:two-leg}. This completes the theorem.

The normal representation for a regular rank-zero orbit also
contains the two-dimensional summand
$\mathfrak u_2\cap E[D]$. It is even-fixed, and
$c_Wf(D)^tc_W^{-1}=f(D)$ makes it odd-fixed as well.
Consequently \cref{m2:tools:passive} applies to this summand and to the
multiplier line in $N_x$.

\subsection{The unitary identity descendant}\label{m2:id:section}\label{m2:sec:identity}

The normal representation for a scalar rank-zero orbit has
an additional moving Lie-algebra summand. Let
\[
 S=G_2\times_\lambda G_2,\qquad
 \mathcal L_{\mathrm{id}}=\Mat_2(E)\oplus\mathfrak u_2,
\]
with action
\begin{equation}\label{m2:id:action}
 (v,h)(Q,Z)=(vQh^{-1},hZh^{-1}).
\end{equation}
For $\mu\in F^\times$, put $A=a_\mu=\diag(1,\mu)$,
and let $D=\zeta I$ with $N_{E/F}(\zeta)=\mu$.
At the scalar representative $A\oplus D$, which is transpose-fixed,
\cref{m2:nor:zero-odd} gives the odd map
\begin{equation}\label{m2:id:odd}
 \Theta(Q,Z)=\bigl(A^{-1}\kappa(Q)\zeta,Z^t\bigr)
             =\bigl(-A^{-1}J\bar QJ\zeta,Z^t\bigr).
\end{equation}
Its left and right factors on $Q$ are $c_U=A^{-1}$ and
$c_W=\zeta^{-1}I$, both of multiplier $\mu^{-1}$.
The map normalizes the even action and is involutive: the two
factors are symmetric, so \eqref{m2:leg:odd-square} gives the identity
on $Q$, while transpose squares to the identity on $Z$.

\begin{theorem}\label{m2:id:vanishing}
For \eqref{m2:id:action} and \eqref{m2:id:odd},
\[
 \Ddist(\Mat_2(E)\oplus\mathfrak u_2)^{S,\Theta=-1}=0.
\]
\end{theorem}

The proof first localizes the nonzero values of a quadratic form on
the moving Lie-algebra factor. Those fibers are the regular two-leg
problems of \cref{m2:leg:regular}. A partial Fourier transform then
places both a distribution and its transform on the same
three-dimensional quadratic null cone. The odd-dimensional
homogeneity theorem finishes the argument.

\subsubsection{The passive central line and the quadratic space}

The equation $Z^*=-Z$ gives
\[
 \mathfrak u_2=
 \left\{\begin{pmatrix}a&b\\c&-\bar a\end{pmatrix}:
          a\in E,\ b,c\in E^-\right\}.
\]
Writing $a=(a-\bar a)/2+(a+\bar a)/2$ yields the direct sum
\begin{equation}\label{m2:id:lie-splitting}
 \mathfrak u_2=E^-I\oplus\mathfrak{su}_{1,1},\qquad
 \mathfrak{su}_{1,1}=
 \left\{\begin{pmatrix}x&b\\c&-x\end{pmatrix}:
          x\in F,\ b,c\in E^-\right\}.
\end{equation}
The central line is fixed pointwise by conjugation and by ordinary
transpose. Projection to it is therefore invariant under the entire
extended action. By \cref{m2:tools:passive}, it suffices to prove the
theorem on $\Mat_2(E)\oplus W$, where
$W=\mathfrak{su}_{1,1}$.

The rational linear map
\begin{equation}\label{m2:id:sl2-map}
 \Phi:W\longrightarrow\mathfrak{sl}_2(F),\qquad
                  \Phi(Z)=P^{-1}ZP
\end{equation}
is a bijection. Indeed, writing $b=\delta y$, $c=\delta z$ gives
\[
 Z=\begin{pmatrix}x&\delta y\\\delta z&-x\end{pmatrix},
 \qquad \Phi(Z)=\begin{pmatrix}x&dy\\z&-x\end{pmatrix}.
\]
All three variables range freely over $F$ and $d\ne0$.
By \cref{m2:simconj:factorization}, every $h\in G_2(F)$ has the form
$h=sPgP^{-1}$, and
\[
 \Phi(hZh^{-1})=g\Phi(Z)g^{-1}.
\]
Conversely every $g\in\GL_2(F)$ occurs through $PgP^{-1}$.
Thus rational conjugation on $W$ is the rational $\GL_2$ adjoint
action, with its scalar kernel. Ordinary transpose
is transported to
\begin{equation}\label{m2:id:transported-transpose}
 X\longmapsto RX^tR^{-1}.
\end{equation}

On $W$ define
\begin{equation}\label{m2:id:quadratic-form}
 B_u(Z,Z')=\tr(ZZ'),\qquad q_u(Z)=\tfrac12\tr(Z^2).
\end{equation}
In the coordinates of \eqref{m2:id:sl2-map},
\[
 q_u(Z)=x^2+dyz=-\det\Phi(Z),
 \qquad
 B_u(Z,Z')=2xx'+d(yz'+zy').
\]
The determinant of the Gram matrix of $B_u$ in coordinates
$(x,y,z)$ is $-2d^2$, which is nonzero in $F$. Thus $W$ is a
nondegenerate quadratic space of $F$-dimension three.
Conjugation preserves $B_u$ by the trace identity, and transpose
does also, since $\tr(Z^tZ'^t)=\tr(Z'Z)=\tr(ZZ')$.
The first $G_2$ factor acts trivially on $W$. It follows that the whole
extended action on $\Mat_2(E)\oplus W$ preserves $q_u$ in the second
variable.

\subsubsection{Nonzero quadratic levels and regular fibers}

For $c\in F^\times$ put
\[
 \mathcal O_c=\{Z\in W:q_u(Z)=c\},\qquad
 X_c=\begin{pmatrix}0&c\\1&0\end{pmatrix},\qquad
 Z_c=PX_cP^{-1}.
\]

\begin{lemma}\label{m2:id:nonzero-orbit}
The level $\mathcal O_c$ is one rational $G_2(F)$-orbit. Its
centralizer $T_c=Z_{G_2}(Z_c)$ is a regular torus, split modulo the
center if $c$ is a square and anisotropic otherwise. Its stabilizer
for the action of $S$ is $S_c=G_2\times_\lambda T_c$.
\end{lemma}

\begin{proof}
Under \eqref{m2:id:sl2-map}, every matrix on this level has
characteristic polynomial $t^2-c$. Since $c\ne0$ and $F$ has
characteristic zero, this polynomial is separable. A trace-zero
matrix with this polynomial is not scalar, so its minimal polynomial
also has degree two. A cyclic vector gives a rational basis in which
the matrix is $X_c$. The rational adjoint action identified above
therefore gives a single $G_2(F)$-orbit.
Its centralizer is
\[
 T_c(F)=E^\times P F[X_c]^\times P^{-1}.
\]
The algebra $F[X_c]\simeq F[t]/(t^2-c)$ is split exactly when $c$ is
a square. Finally $S$ acts on the base through its right factor $h$,
so fixing $Z_c$ means $h\in T_c$, with $v$ subject only to the
common-multiplier condition. This is exactly $S_c$.
\end{proof}

The projection
$\Mat_2(E)\times\mathcal O_c\longrightarrow\mathcal O_c$
is equivariant, with fiber $\Mat_2(E)$ at $Z_c$ and stabilizer
$S_c$. Both $S$ and $S_c$ are reductive and unimodular, so Frobenius
descent carries the even invariant distributions to ordinary
$S_c$-invariant distributions on this fiber, with no modular twist.
To use it for the odd character we now compute the transporter and
the induced odd operator explicitly.

From \eqref{m2:id:transported-transpose},
\[
 RX_c^tR^{-1}=\begin{pmatrix}0&d\\c/d&0\end{pmatrix}.
\]
The rational diagonal matrix $g_c=\diag(c,d)$ conjugates this
matrix to $X_c$. Therefore
\begin{equation}\label{m2:id:level-transporter}
 h_c=Pg_cP^{-1}\in G_2(F),\qquad
 h_cZ_c^th_c^{-1}=Z_c,\qquad \lambda(h_c)=cd.
\end{equation}
Choose $v_c\in G_2(F)$ with $\lambda(v_c)=cd$; one explicit choice
is $v_c=h_c$. The pair $(v_c,h_c)\in S(F)$ corrects the base
action of $\Theta$ at $Z_c$. On the $Q$-fiber its odd action is
\begin{equation}\label{m2:id:fiber-odd}
 Q\longmapsto v_cA^{-1}\kappa(Q)\zeta h_c^{-1}
                 =c_U'\kappa(Q)(c_W')^{-1},
 \qquad c_U'=v_cA^{-1},\quad c_W'=h_c\zeta^{-1}.
\end{equation}
These factors satisfy
\[
 \lambda(c_U')=\lambda(c_W')=cd/\mu,
 \qquad c_W'Z_c^t(c_W')^{-1}=Z_c.
\]
To apply \cref{m2:leg:regular}, observe that $Z_c^2=cI$ and $Z_c^*=-Z_c$,
so
\[
 Z_c^*Z_c=-cI.
\]
Thus $Z_c$ itself is an invertible regular semisimple element of
$G_2(F)$, of multiplier $-c$. Taking $D=Z_c$ in that theorem gives
the torus $T_c$, the even group $S_c$, and the odd map
\eqref{m2:id:fiber-odd}.

The homogeneous base has a positive $S$-invariant measure because
$S$ and $S_c$ are unimodular. Transpose preserves its orbit and
normalizes the action; its positive measure-scaling constant is
one because its square is the identity. The correcting even
element also preserves the base measure. Naturality of the
Frobenius construction with this measure therefore intertwines
the odd action on the total space with the corrected fiber action
\eqref{m2:id:fiber-odd}.
Thus \cref{m1:prop:odd-Frobenius} and \cref{m2:leg:regular} prove
\begin{equation}\label{m2:id:nonzero-vanishing}
 \Ddist(\Mat_2(E)\times\mathcal O_c)^{S,\Theta=-1}=0
                         \qquad(c\ne0).
\end{equation}
The sign is trivial on the even fiber stabilizer and is minus one
on its corrected odd coset.

\subsubsection{Support and partial Fourier transform}

Let
\[
 T\in\Ddist(\Mat_2(E)\oplus W)^{S,\Theta=-1}.
\]
On the invariant open set $q_u(Z)\ne0$, the continuous map
$(Q,Z)\mapsto q_u(Z)\in F^\times$ has the complete fibers
treated in \eqref{m2:id:nonzero-vanishing}. Bernstein localization
therefore makes the restriction of $T$ to this open set zero.
The supported-distribution exact sequence gives
\begin{equation}\label{m2:id:support}
 \supp(T)\subset\Mat_2(E)\times\{Z\in W:q_u(Z)=0\}.
\end{equation}

Choose a nontrivial additive character $\psi$ of $F$ and the
self-dual Haar measure for $B_u$. Let $\mathcal F_W$ be partial
Fourier transform in the $W$-variable, using the kernel
$\psi(B_u(Z,Z'))$. Every even or odd transformation on $W$ is an
isometry of $B_u$. If its matrix is $L$ and the Gram matrix of the
pairing is $H$, the identity $L^tHL=H$ implies
$|\det L|_F=1$. Changing variables in the Fourier integral
therefore introduces no absolute determinant factor and gives
\[
 \mathcal F_W L_*=L_*\mathcal F_W.
\]
The same statement holds on distributions by duality. The
$Q$-transformations in \eqref{m2:id:action} and \eqref{m2:id:odd} are
linear in $Q$ and independent of $Z$; they commute with this partial
transform. Hence
\[
 \widehat T=\mathcal F_WT
       \in\Ddist(\Mat_2(E)\oplus W)^{S,\Theta=-1}.
\]
Applying the same open localization to $\widehat T$ gives the
second support statement
\begin{equation}\label{m2:id:fourier-support}
 \supp(\widehat T)\subset
              \Mat_2(E)\times\{Z\in W:q_u(Z)=0\}.
\end{equation}
We can now invoke the non-Archimedean homogeneity theorem
\cref{m2:tools:homogeneity} (\cite[Theorem~5.1.5]{AGDescent})
with parameter manifold $\Mat_2(E)$ and nondegenerate quadratic
space $(W,B_u)$. Its quadratic null cone is
$\{Z:B_u(Z,Z)=0\}=\{Z:q_u(Z)=0\}$, since $2\ne0$ in $F$.
Its two hypotheses are exactly
\eqref{m2:id:support} and \eqref{m2:id:fourier-support}, using partial
Fourier transform in $W$. The final clause of that theorem gives
zero whenever $\dim_F W$ is odd. Here this dimension is three, so
$T=0$. Restoring the passive line in \eqref{m2:id:lie-splitting}
proves \cref{m2:id:vanishing}.

\begin{remark}\label{m2:id:nilpotent-remark}
The zero level of $q_u$ consists of the origin and the single nonzero
rational nilpotent adjoint orbit. The latter assertion follows from
the rational $\GL_2$ adjoint description in \eqref{m2:id:sl2-map}.
The support inclusions and homogeneity apply to this entire level.
\end{remark}

\subsection{Proof of the unitary multiplicity-one theorem}
\label{m2:sec:completion}

\subsubsection{Invariance under the corrected involutions}

The nine normal representations, with their rational stabilizers,
odd operators, and parameter domains, are listed in
\cref{m2:nor:classification}. Their vanishing results are:
\begin{itemize}
\item \cref{m2:id:vanishing} for the scalar rank-zero family;
\item \cref{m2:leg:regular} for both regular rank-zero families;
\item \cref{m2:simconj:vanishing} for the dense scalar family;
\item \cref{m2:elem:vanishing} for the other dense families and both
rank-one families.
\end{itemize}

\begin{proposition}\label{m2:completion:normal-vanishing}
For every closed rational $K(F)$-orbit in $G_4(F)$, with
representative $x$ and extended stabilizer $\Kt_x(F)$,
\[
 \Ddist(N_x(F))^{\Kt_x(F),\eps}=0.
\]
\end{proposition}

\begin{proof}
By \cref{m2:rat:classification,m2:rat:transpose}, the nine rational
families are exhaustive and every extended stabilizer has an odd
coset. Its normal action is the one in \cref{m2:nor:classification}.
In the two regular rank-zero rows, the transporter factors have
equal multipliers and satisfy $c_WD^tc_W^{-1}=D$. Thus
\cref{m2:leg:regular} applies to the $M_2(E)$ summand with its
corrected odd action. The remaining $F^2$ is fixed by both actions,
so \cref{m2:tools:passive} adjoins it. For the scalar rank-zero
row, \cref{m2:id:vanishing} applies to the entire reduced normal space
$M_2(E)\oplus\mathfrak u_2$. The remaining six rows are covered by
\cref{m2:elem:vanishing,m2:simconj:vanishing}.
Finally \cref{m2:nor:full} gives
$N_x=N_x^0\oplus F[I_4/2]$, with the last line fixed by the
extended stabilizer. Applying \cref{m2:tools:passive} proves the assertion.
\end{proof}

\subsubsection{Generalized descent for the residual action}

\begin{theorem}\label{m2:completion:residual}
Every $K(F)$-invariant distribution on $G_4(F)$ is fixed by
ordinary transpose. Equivalently,
\[
 \Ddist(G_4(F))^{\Kt(F),\eps}=0.
\]
\end{theorem}

\begin{proof}
We apply \cref{m1:cor:HC-invariance} to the action of $K$ on
$G_4$, taking $\tau$ to be ordinary transpose. The group $K$ is connected
reductive, and $G_4$ is smooth affine. The action
\eqref{m2:red:action}, transpose, and \eqref{m2:red:alpha} are
algebraic over $F$. By \cref{m2:rat:transpose}, every closed rational
orbit is transpose-stable, and its extended stabilizer is
$K_x(F)\sqcup K_x(F)k_x\sigma$. Its normal space and sign action
are those of \cref{m2:nor:full};
\cref{m2:completion:normal-vanishing} supplies the required vanishing.
The corollary therefore gives transpose invariance, equivalently
the required vanishing of sign-equivariant distributions.
\end{proof}

\begin{corollary}\label{m2:completion:biH}
Every bi-$H_{\mathrm U}$-invariant distribution on $G_{\mathrm U}$ is fixed by
$\tau(g,z)=(g^t,z^t)$.
\end{corollary}

\begin{proof}
The modular factor and compatibility with the invariant base measure
for the isomorphism \eqref{m2:red:distiso} were verified in
\cref{m2:red:frobenius}. Its image is fixed by fiber transpose
by \cref{m2:completion:residual}, so its inverse image is fixed by
$\tau$.
\end{proof}

\subsubsection{The product inequality and equality of multiplicities}

\begin{proof}[Proof of \cref{thm:main-unitary}]
Fix an irreducible admissible smooth complex representation
$\pi$ of $G_{\mathrm U}$, and put
$m(\pi)=\dim\Hom_{H_{\mathrm U}}(\pi,\mathbb C)$.
The groups $G_{\mathrm U}$ and $H_{\mathrm U}$ are unimodular, so $G_{\mathrm U}/H_{\mathrm U}$
has an invariant measure. The anti-involution preserves the
subgroup by \cref{m2:red:transpose}, and fixes bi-invariant
distributions by \cref{m2:completion:biH}. Thus the
non-Archimedean Gelfand--Kazhdan criterion
\cite[Lemma~4.2]{Prasad90}, in the form of \cref{m2:tools:gk},
gives
\[
 m(\pi)m(\pi^\vee)\leq1.
\]
Finally \cref{m2:red:duality}, obtained from
Roche--Vinroot's dualizing theorem and the preservation of $H_{\mathrm U}$
by inverse transpose, gives $m(\pi)=m(\pi^\vee)$. Their common
dimension is therefore at most one.
\end{proof}
